\chapter{The Fourier Transform}\label{ch:the_Fourier_transform}
  Let \(f(\mathbf{x})\) be absolutely integrable on \(\R^{n}\). Its \textbf{Fourier transform}\index{Fourier transform} \((\mc{F}f)(\boldsymbol{\xi})\) is defined by
  \[
    (\mc{F}f)(\boldsymbol{\xi}) = \int_{\R^{n}}f(\mathbf{x})e^{-2\pi i\<\boldsymbol{\xi},\mathbf{x}\>}\,d\mathbf{x},
  \]
  for \(\boldsymbol{\xi} \in \R^{n}\). It is absolutely convergent precisely because \(f(\mathbf{x})\) is absolutely integrable on \(\R^{n}\). Let us first demonstrate some properties of the Fourier transform.

  \begin{proposition}
    Let \(f(\mathbf{x})\) and \(g(\mathbf{x})\) be absolutely integrable on \(\R^{n}\). Then the following properties hold:
    \begin{enumerate}[label*=(\roman*)]
      \item For any \(\a,\b \in \R\), we have
      \[
        (\mc{F}(\a f+\b g))(\boldsymbol{\xi}) = \a(\mc{F}f)(\boldsymbol{\xi})+\b(\mc{F}g)(\boldsymbol{\xi}).
      \]
      \item If \(g(\mathbf{x}) = f(\mathbf{x}+\boldsymbol{\a})\) for any \(\boldsymbol{\a} \in \R^{n}\) then
      \[
        (\mc{F}g)(\boldsymbol{\xi}) = e^{2\pi i\<\boldsymbol{\a},\boldsymbol{\xi}\>}(\mc{F}f)(\boldsymbol{\xi}).
      \]
      \item If \(g(\mathbf{x}) = f(A\mathbf{x})\) for any \(A \in \GL_{n}(\R)\) then
      \[
        (\mc{F}g)(\boldsymbol{\xi}) = \frac{1}{|\det(A)|}(\mc{F}f)((A^{-1})^{t}\boldsymbol{\xi}).
      \]
    \end{enumerate}
  \end{proposition}
  \begin{proof}
    Property (i) is immediate from linearity of the integral while property (ii) follows by applying the change of variables \(\mathbf{x} \mapsto \mathbf{x}-\boldsymbol{\a}\). Property (iii) is proved by performing the change of variables \(\mathbf{x} \mapsto A^{-1}\mathbf{x}\) which has Jacobian matrix \(A^{-1}\) with Jacobian determinant \(\frac{1}{|\det(A)|}\).
  \end{proof}

  The Fourier transform is intimately related to periodic functions. Let \(\L\) be a complete integral lattice in \(\R^{n}\) with fundamental domain \(\mc{M}\) and denote the dual integral lattice by \(\L^{\vee}\). Suppose \(f(\mathbf{x})\) is \(\L\)-periodic and integrable on \(\mc{M}\). Then we define the \(\l^{\vee}\)-th \textbf{Fourier coefficient}\index{Fourier coefficient} \(\hat{f}(\l^{\vee})\) of \(f(\mathbf{x})\) by
  \[
    \hat{f}(\l^{\vee}) = \int_{\mc{M}}f(\mathbf{x})e^{-2\pi i\<\l^{\vee},\mathbf{x}\>}\,d\mathbf{x}.
  \]
  The \textbf{Fourier series}\index{Fourier series} \(S_{f}(x)\) is defined by
  \[
    S_{f}(\mathbf{x}) = \frac{1}{V_{\L}}\sum_{\l^{\vee} \in \L^{\vee}}\hat{f}(\l^{\vee})e^{2\pi i\<\l^{\vee},\mathbf{x}\>}.
  \]
  Notice that the Fourier coefficients are only defined for elements of the dual integral lattice and that the Fourier series is a sum over these elements. A crucial question is where the Fourier series of \(f(\mathbf{x})\) converges, if at all, and if so does it even converge to \(f(\mathbf{x})\) itself. Under quite reasonable conditions the Fourier series converges uniformly to the function itself, but we first make a reduction. Fix a basis \(\l_{1},\ldots,\l_{n}\) for \(\L\) and let \(P\) be the associated generator matrix. This means \(P\) is the base change matrix from the standard basis to \(\l_{1},\ldots,\l_{n}\) so that \(P\) is the unique invertible linear transformation on \(\R^{n}\) satisfying \(\L = P\Z^{n}\) and \(\L^{\vee} = (P^{-1})^{t}\Z^{n}\). Letting \( f_{P}(\mathbf{x}) = f(P\mathbf{x})\), we see that \(f_{P}(\mathbf{x})\) is \(1\)-periodic in each component and integrable on \([0,1]^{n}\). As \(\Z^{n}\) is self-dual and its covolume is \(1\), the \(\mathbf{n}\)-th Fourier coefficient \(\hat{f}_{P}(\mathbf{n})\) of \(f_{P}(\mathbf{x})\) is given by
  \[
    \hat{f}_{P}(\mathbf{n}) = \int_{[0,1]^{n}}f(\mathbf{x})e^{-2\pi i\<\mathbf{n},\mathbf{x}\>}\,d\mathbf{x}.
  \]
  and the Fourier series is
  \[
    S_{f_{P}}(\mathbf{x}) = \sum_{\mathbf{n} \in \Z^{n}}\hat{f}_{P}(\mathbf{n})e^{2\pi i\<\mathbf{n},\mathbf{x}\>}.
  \]
  As \(P\) is invertible, it suffices to study the convergence properties of this Fourier series. The reduction here this that we may assume our Fourier series are \(1\)-periodic in each component and in this case we have the following well-known result:

  \begin{theorem}
    Suppose \(f(\mathbf{x})\) is a smooth function on \(\R^{n}\) and is \(1\)-periodic in each component. Then the Fourier series of \(f(\mathbf{x})\) converges uniformly everywhere to \(f(\mathbf{x})\).
  \end{theorem}

  In the case of \(1\)-periodic functions of a single variable, we can do better as we may merely assume \(f(x)\) is of bounded variation. This is known as the \textbf{Dirichlet-Jordan test}\index{Dirichlet-Jordan test}.

  \begin{theorem*}[Dirichlet-Jodan test]
    Suppose \(f(x)\) is a function on \(\R\) which is \(1\)-periodic and of bounded variation. Then the Fourier series of \(f(x)\) converges locally uniformly to \(f(x)\) on every set where \(f(x)\) is continuous. Moreover, at any jump discontinuity, the Fourier series of \(f(x)\) converges to the average of the left-hand and right-hand limits of \(f(x)\). In particular, this holds for all continuously differentiable functions with at most a finite number of jump discontinuities.
  \end{theorem*}